\documentclass[11pt]{article}

% ---------------------------------------------------------------------------
%  L4 — The by-height generating function is not D-finite.
%  Category L (entirely machine-written); see paper/README.md and
%  docs/publication-split.md.
%
%  Source material, all in this repository:
%    results/anisotropic-not-dfinite.md   the theorem, the boxes, the prior art
%    results/diagonal-formula.md        the Atom Ledger and root separation
%    results/growth-constant.md     the crude lambda bound the proof uses
%    results/fixed_height_gfs.txt         the banked G_H
%    papers/MISSING.md                    the literature searches, per database
% ---------------------------------------------------------------------------

% ---------------------------------------------------------------------------
%  paper/shared/preamble.tex — packages and macros common to every manuscript
%  in this directory.  Factored out of paper/polyplets-report.tex and
%  paper/technical-report.tex, which between them had two copies of the same
%  twenty lines.
%
%  technical-report.tex does NOT \input this file and is not to be edited to do
%  so: it is jasonp's prose and is read-only to the machine (paper/README.md,
%  "Who may edit what").  The duplication there is deliberate.
%
%  Assumes \documentclass[11pt]{article} has already been issued.
% ---------------------------------------------------------------------------
\usepackage[margin=1in]{geometry}
\usepackage{amsmath,amssymb}
\usepackage{amsthm}
\usepackage{booktabs}
\usepackage{graphicx}
\usepackage{numprint}
\usepackage{microtype}
\usepackage[table]{xcolor}
\usepackage[hidelinks]{hyperref}
\usepackage{flafter}   % a float never appears before the text that introduces it
\usepackage{float}     % [H] pins tables/figures exactly where written
\usepackage[font=small,labelfont=bf,labelsep=period]{caption}
\setlength{\intextsep}{16pt plus 3pt minus 3pt}

\npthousandsep{,}
\npfourdigitsep
\newcommand{\num}[1]{\numprint{#1}}
\newcommand{\g}{\cellcolor{gray!15}}

% Theorem environments.  One counter shared across kinds, so that a reference
% like "Theorem 4" is unambiguous in a paper that also has lemmas and
% propositions.
\theoremstyle{plain}
\newtheorem{theorem}{Theorem}
\newtheorem{proposition}[theorem]{Proposition}
\newtheorem{lemma}[theorem]{Lemma}
\newtheorem{corollary}[theorem]{Corollary}
\theoremstyle{definition}
\newtheorem{definition}[theorem]{Definition}
\newtheorem{example}[theorem]{Example}
\theoremstyle{remark}
\newtheorem{remark}[theorem]{Remark}
\newtheorem{conjecture}[theorem]{Conjecture}
\newtheorem{openproblem}[theorem]{Open Problem}

% Repo-path citations.  Every one of these must resolve in the tree that ships
% with the paper; `make gate-citations` checks the markdown, and
% paper/verify_claims.py checks the manuscripts.
\newcommand{\repo}[1]{\texttt{#1}}

% OEIS links.
\newcommand{\oeis}[1]{\href{https://oeis.org/#1}{#1}}

% Growth constants, so a rename is one edit.
\newcommand{\lam}{\lambda}
\newcommand{\muH}{\mu_H}

% ---------------------------------------------------------------------------
%  paper/shared/disclosure.tex — the two authorship disclosure blocks.
%
%  docs/publication-split.md §1 fixes the rule these implement:
%
%    "Under no circumstances will there be any ambiguity about whether a paper
%     is wholly my work (none), merely my prose (some) or entirely LLM-authored
%     (some)."   — jasonp, 2026-08-07
%
%  There are exactly two blocks and they live in exactly one file, so that a
%  paper cannot quietly soften its own disclosure.  The fixed sentences take no
%  arguments.  The ONE thing a paper supplies is the per-result verification
%  ledger, because that is the part that differs honestly between papers — and
%  §1 requires it to be per result, not in aggregate.
%
%  Do not add a \Pdisclosure or \Ldisclosure variant with an optional softening
%  argument.  If a paper does not fit one of these two blocks, the paper is
%  wrong, not the block.
% ---------------------------------------------------------------------------

\newsavebox{\disclosurebox}

\newenvironment{disclosureframe}
  {\begin{center}\begin{lrbox}{\disclosurebox}\begin{minipage}{0.93\textwidth}\small}
  {\end{minipage}\end{lrbox}\fbox{\usebox{\disclosurebox}}\end{center}}

% --- P: jasonp's prose -------------------------------------------------------
% Byline: Jason Parker, alone.  Argument: the per-result ledger of what the
% machine did.
\newcommand{\Pdisclosure}[1]{%
\begin{disclosureframe}
\textbf{Authorship.}\enspace Every sentence of this paper was written by its
named author. He has read every proof it states and believes each one, and he
is responsible for the correctness of what it claims.

\smallskip
\textbf{Machine contribution.}\enspace #1
\end{disclosureframe}}

% --- L: entirely machine-written --------------------------------------------
% Argument: the per-result verification ledger — for each result, what warrants
% it (Lean kernel, exact certificate, or labelled measurement) and how much of
% it a human has checked.
\newcommand{\Ldisclosure}[1]{%
\begin{disclosureframe}
\textbf{Authorship.}\enspace This document was written end to end by a large
language model. That includes its mathematics, its prose, its tables and its
\LaTeX{}. No sentence of it was written by a human. The mathematics it reports
was likewise derived by machine.

\smallskip
\textbf{Human role.}\enspace The person who directed this work chose what to
compute, chose what to publish, and is responsible for the decision to release
this document. Except where the ledger below says otherwise, that person has
not personally checked the arguments here, and this disclosure is not to be
read as an endorsement of them.

\smallskip
\textbf{What warrants each result.}\enspace #1
\end{disclosureframe}}

% --- Draft banner ------------------------------------------------------------
% For an L-paper whose verification ledger is not yet settled.  Loud on purpose:
% an L-paper without a truthful ledger must not be mistaken for a finished one.
%
% \firstdraftbanner is the standard wording, which five papers previously
% carried character-for-character in their own preambles.  It lives here for the
% same reason the disclosure blocks do: identical text in seven places drifts,
% and a paper must not be able to soften it by editing its own copy.  The
% optional argument is for a gate specific to one paper, appended verbatim; a
% paper needing different wording (L6, compute-gated) calls \draftbanner direct.
\newcommand{\firstdraftbanner}[1][]{%
\draftbanner{This is a first draft. The verification ledger below records what
warrants each result \emph{mechanically}; the column that records what a human
has read is empty, deliberately, because nobody has read it yet. Do not
circulate until that column says something true.#1}}

\newcommand{\draftbanner}[1]{%
\begin{center}
\fcolorbox{red!60!black}{yellow!15}{%
  \begin{minipage}{0.93\textwidth}\small
  \textbf{DRAFT — NOT FOR CIRCULATION.}\enspace #1
  \end{minipage}}
\end{center}}


\newcommand{\Z}{\mathbb{Z}}
\newcommand{\Q}{\mathbb{Q}}
\newcommand{\C}{\mathbb{C}}
\newcommand{\F}{\mathbb{F}}

\title{\textbf{An arithmetic obstruction to D-finiteness for
lattice-animal height generating functions}\\[4pt]
\large Bounded degree and bounded house, against a strictly
increasing family of slice growth constants}
\author{[byline: jasonp's call --- \repo{docs/publication-split.md} \S1]}
\date{August 2026}

\begin{document}
\maketitle

\firstdraftbanner

\Ldisclosure{%
\emph{Theorem~\ref{thm:main} (not D-finite)} --- proved below from four
classical ingredients (Fatou, Pringsheim, Perron--Frobenius, Northcott) plus
Lemma~\ref{lem:mono}. Not formalized. Its opening move --- extracting a
coefficient relation from the ODE --- is not ours and is cited to
Bousquet-M\'elou and Rechnitzer~\cite[Lemma~9]{bmr2002}.
\emph{Lemma~\ref{lem:mono} (strict monotonicity)} --- proved below by
Perron--Frobenius; the numerical values in Table~\ref{tab:mu} are labeled
measurement, and the strictness the proof needs does not depend on them.
\emph{The exclusion boxes of \S\ref{sec:effective}} --- rest on
certificates computed modulo a $61$-bit prime, where a gcd equal to $1$ with
preserved degrees is a rigorous certificate over $\Q$; the irreducibility
certificates for $H \le 8$ are single-prime for $H \le 4$ and multi-prime
subset-sum-intersection for $H = 5..8$.
\emph{The literature comparison of \S\ref{sec:related}} --- three databases,
one of them full-text, recorded per search in \repo{papers/MISSING.md}. No
collision found, and no claim here rests on that.
\emph{Human verification:} none, as of this draft.}

\begin{abstract}
\noindent
Let $T(n,H)$ count the fixed king-lattice animals of $n$ cells whose bounding
box has height exactly $H$, and let
$F(x,y) = \sum_{n,H} T(n,H)x^ny^H$. We prove that $F$ is not D-finite.

The obstruction is arithmetic rather than topological or analytic, which is
what makes the proof cheap to transport. Write $\mu_H$ for the growth constant
of the height-$\le H$ strip. The proof needs two facts about the family
$(\mu_H)$: it is \emph{strictly increasing} (Perron--Frobenius, since each strip system sits
inside the next as a proper principal submatrix of an irreducible nonnegative
matrix), and it is \emph{bounded} (by $\lambda \le 9.3154$, or by any crude
bound). Each $\mu_H$ is an algebraic integer all of whose conjugates are at most
$\mu_H$ in modulus, so Northcott finiteness forces the algebraic degrees
$[\Q(\mu_H):\Q]$ to be unbounded. Against that, a D-finite $F$ would force those
degrees to be bounded by the $x$-degree of its annihilating operator at all but
finitely many heights. Contradiction.

The proof uses only integer counts, rationality of each height slice, and a
strictly increasing bounded family of slice growth constants. It therefore
applies verbatim to fixed polyominoes by height, to polyhexes, and to any
lattice-animal family with a column transfer matrix.

The argument is also effective. From certified data we exclude explicit boxes:
no annihilating $y$-ODE exists of
order $r \le 5$ with $x$-degrees $\le 28$, of order $r \le 3$ with $x$-degrees
$\le 180$, of order $r \le 1$ with $x$-degrees $\le 1253$, and so on.

The method has a ceiling. The step ``a rational polynomial of degree $\le D$
cannot vanish at an algebraic number of degree $> D$'' is the whole argument, and
it fails the instant the coefficients are allowed to be analytic rather than
polynomial. So the result is not-D-finite over $\Q(x)$ and does \emph{not}
extend to Klazar's larger $D_A$-finite class.
\end{abstract}

\tableofcontents

\section{Introduction}

Counting lattice animals by size alone gives a sequence, $a(n)$, about which
almost nothing exact is known. Counting them by size \emph{and height} gives a
two-variable object, and the extra variable turns out to be what makes a proof
possible. This is Rechnitzer's observation in the haruspicy
program~\cite{rechnitzer2003,rechnitzer2006}: the anisotropic generating
function can be shown not to be D-finite even when the isotropic one is out of
reach, and the mechanism is that the height slices are individually rational and
their poles proliferate.

We prove the corresponding statement for the king lattice, by a different
mechanism.

\begin{definition}
$T(n,H)$ is the number of fixed polyplets (animals on the king lattice)
of $n$ cells whose
bounding box has height exactly $H$; $G_H(x) := \sum_n T(n,H)x^n$; and
\[
  F(x,y) \;:=\; \sum_{n,H} T(n,H)\,x^n y^H \;=\; \sum_H G_H(x)\,y^H .
\]
\end{definition}

\begin{theorem}
\label{thm:main}
$F(x,y)$ is not D-finite.
\end{theorem}

\subsection{The shape of the argument}

Every proof of non-D-finiteness in this area starts the same way. If $F$ satisfies
$\sum_{i \le r}p_i(x,y)\,\partial_y^i F = 0$ with $p_i \in \Q[x,y]$ not all zero,
extracting the coefficient of $y^H$ turns the ODE into a linear recurrence in $H$
whose coefficients are polynomials in $x$ and $H$. That step is due to
Bousquet-M\'elou and Rechnitzer~\cite[Lemma~9]{bmr2002} and we cite it rather than re-deriving it.

What is done with the recurrence is where the routes diverge, and there are three
of them:

\begin{table}[H]
\centering\small
\begin{tabular}{l l l}
\toprule
 & obstruction & lever \\
\midrule
Bousquet-M\'elou--Rechnitzer 2002~\cite{bmr2002} & topological & accumulation of the pole set in $\C$ \\
Klazar 2003~\cite{klazar2003} & analytic & zero radius of convergence vs.\ forced analyticity \\
this paper & \textbf{arithmetic} & degree and house of algebraic numbers \\
\bottomrule
\end{tabular}
\caption{Three obstruction types in the same neighborhood.}
\label{tab:three}
\end{table}

\section{Ingredients}

\subsection{Rationality, and \texorpdfstring{$\mu_H$}{mu\_H} as an algebraic integer}

\begin{lemma}
\label{lem:alg}
Each $G_H$ is rational with integer coefficients; $\mu_H$, the reciprocal of its
radius of convergence, is an algebraic integer; and every Galois conjugate of
$\mu_H$ has modulus at most $\mu_H$, so $\mathrm{house}(\mu_H) = \mu_H$.
\end{lemma}

\begin{proof}
$G_H$ is the generating function of a finite transfer matrix, hence rational, and
its coefficients are counts. By Fatou's lemma~\cite{fatou1906}, a rational power
series with integer coefficients has a lowest-terms denominator $Q_H$ with
$Q_H(0) = 1$ and integer coefficients. The reversal of $Q_H$ is then monic and
integral, and its roots are exactly the reciprocals of the poles of $G_H$.

Since $T \ge 0$, Pringsheim's theorem~\cite[Thm.~IV.6]{flajoletSedgewick2009}
puts a singularity of $G_H$ at the positive-real point of its circle of
convergence; for a rational function that singularity is a pole. So $\mu_H$ is
the reciprocal of a pole, hence a root of the monic integral reversal, hence an
algebraic integer.

Every pole of $G_H$ has modulus at least the radius $1/\mu_H$, by definition of
the radius of convergence. So every root of the reversal, and in particular
every conjugate of $\mu_H$, has modulus at most $\mu_H$.
\end{proof}

The house bound is then $\mathrm{house}(\mu_H) = \mu_H < \lambda \le 9.3154$,
the certified upper bound of companion paper L3, \emph{A certified two-sided
bound for the polyplet growth constant}. Any crude bound would do: the
elementary $\lambda \le 5^5/4^4 \approx 12.2$ suffices.

\subsection{Strict monotonicity}

\begin{lemma}[strict monotonicity]
\label{lem:mono}
$\mu_{H+1} > \mu_H$ for every $H$. Consequently the $\mu_H$ are pairwise
distinct, and $1/\mu_H$ is a pole of $G_H$ that is not a pole of any $G_j$ with
$j < H$.
\end{lemma}

\begin{proof}
Let $\nu_H$ be the growth constant of the height-$\le H$ strip counts, computed
by the column-signature transfer matrix $M_H(x)$: states are boundary signatures
under the nonempty-column invariant. The signature digraph is strongly connected:
any state reaches the single-cell state by closing off its components with a
spanning column, and any state is reachable from the single-cell state. So
$M_H(x)$ is irreducible and nonnegative for $x > 0$.

$M_H$ sits inside $M_{H+1}$ as a proper principal submatrix, namely the
signatures that do not use row $H+1$, and the extra states connect back through
the explicit cycle: single cell $\to$ full $(H+1)$-column $\to$ single cell.
Deleting rows and columns of an irreducible nonnegative matrix strictly decreases
the Perron root~\cite[Ch.~8]{hornJohnson2013}, so $\rho_{H+1}(x) > \rho_H(x)$ at
every $x > 0$. At $x^*_H$, where $\rho_H = 1$, this gives
$\rho_{H+1}(x^*_H) > 1$, hence $x^*_{H+1} < x^*_H$ and $\nu_{H+1} > \nu_H$
strictly.

Exact-height counts are second differences of strip counts, and strictness rules
out cancellation there, so $\mu_H = \nu_H$ and the pole $1/\mu_H$ of $G_H$
is genuine.
\end{proof}

\begin{table}[H]
\centering\small
\begin{tabular}{r r r r r r r r r r}
\toprule
$H$ & 1 & 2 & 3 & 4 & 5 & 6 & 7 & 8 & 9 \\
\midrule
$\mu_H$ & 1.0 & 2.41421 & 3.44372 & 4.18232 & 4.71780 & 5.11532 & 5.41785 & 5.65337 & 5.84046 \\
\bottomrule
\end{tabular}
\caption{Strip growth constants, by exact-coefficient evaluation
($\mu_{10} = 5.99170$). Labelled measurement; Lemma~\ref{lem:mono} does not
depend on these values, and the proof of Theorem~\ref{thm:main} does not use
them. $\mu_2 = 1+\sqrt2$ exactly.}
\label{tab:mu}
\end{table}

\subsection{Northcott finiteness}

\begin{lemma}
\label{lem:northcott}
$[\Q(\mu_H):\Q] \to \infty$.
\end{lemma}

\begin{proof}
Fix $D$. An algebraic integer of degree $\le D$ whose conjugates all have modulus
at most $B$ has minimal polynomial with integer coefficients bounded by
$\binom{D}{k}B^k$ in absolute value --- they are elementary symmetric functions
of the conjugates. There are therefore finitely many such integer polynomials,
hence finitely many such algebraic numbers~\cite{northcott1949,kronecker1857}.

By Lemma~\ref{lem:alg} every $\mu_H$ is an algebraic integer of house at most
$B = 9.3154$, and by Lemma~\ref{lem:mono} they are pairwise distinct, hence
infinitely many. So for every $D$ only finitely many have degree $\le D$.
\end{proof}

\section{The dichotomy}

\begin{theorem}[dichotomy]
\label{thm:dichotomy}
Suppose $F$ is annihilated by a $y$-ODE of order $r$ with $x$-degrees at most
$D$. Then $[\Q(\mu_H):\Q] \le D$ for all but at most $r$ values of $H \ge r$.
\end{theorem}

\begin{proof}
Write the annihilating relation as $\sum_{i \le r}p_i(x,y)\partial_y^iF = 0$ with
$p_i = \sum_j p_{ij}(x)y^j$ in $\Q[x,y]$, not all zero, and $D = \max_i \deg_x p_i$.
Write $(N)_{(i)} := N(N-1)\cdots(N-i+1)$ for the falling factorial. Since
$\partial_y^iF = \sum_k (k+i)_{(i)}G_{k+i}\,y^k$, the $(i,j)$ term
contributes $p_{ij}(x)\,(H+i-j)_{(i)}\,G_{H+i-j}$ to the coefficient of $y^H$, so
extracting $y$-coefficients~\cite[Lemma~9]{bmr2002} gives, for every $H \ge 0$,
\begin{equation}
  \sum_{i,j} p_{ij}(x)\,(H+i-j)_{(i)}\,G_{H+i-j}(x) \;=\; 0 .
  \label{eq:rec}
\end{equation}

The height indices here are \emph{not} bounded above by $H$: the $(i,j)$ term
sits at height $H + i - j$, which exceeds $H$ whenever $i > j$. For
$\partial_yF = 0$ the only term is $G_{H+1}$ and no term of height $H$ occurs at
all. So the face to run the argument on is the top one. Put
\[
  m \;:=\; \max\{\,i - j \;:\; p_{ij} \not\equiv 0\,\} \;\le\; r,
  \qquad
  c(x,H) \;:=\; \sum_{i} p_{i,\,i-m}(x)\,(H+m)_{(i)} ,
\]
so that $c$ is the coefficient of the highest slice $G_{H+m}$ in
\eqref{eq:rec}. Then $\deg_x c \le D$ and $\deg_H c \le r$, and
$c \not\equiv 0$: the falling factorials $(H+m)_{(i)}$ have distinct degrees in
$H$, hence are linearly independent, and $m$ was chosen so that at least one
$p_{i,i-m}$ survives.

Fix $H_0 \ge r$ and instance \eqref{eq:rec} at $H = H_0 - m$, which is a legal
index since $m \le r$. The top slice is then $G_{H_0}$, with coefficient
$c(x,H_0-m)$, and every other term refers to a height strictly below $H_0$.
Evaluate near $x = 1/\mu_{H_0}$: by Lemma~\ref{lem:mono} the lower slices have
strictly larger radius of convergence and so are finite there, while $G_{H_0}$
has a genuine pole. Hence $c(1/\mu_{H_0},H_0-m) = 0$.

Now $c(\cdot,H_0-m)$ is a polynomial over $\Q$ of degree at most $D$. If
$[\Q(\mu_{H_0}):\Q] > D$, then $1/\mu_{H_0}$ also has degree $> D$ over $\Q$ and
no nonzero rational polynomial of degree $\le D$ vanishes at it; so
$c(\cdot,H_0-m) \equiv 0$.

Suppose that happens at $r+1$ distinct heights $H_0 \ge r$. The polynomial
$c(x,\cdot)$ has degree at most $r$ in its second argument and vanishes at the
$r+1$ distinct points $H_0 - m$, for each fixed $x$; so $c \equiv 0$,
contradicting $c \not\equiv 0$.
\end{proof}

\begin{proof}[Proof of Theorem~\ref{thm:main}]
Suppose $F$ were D-finite. Without loss of generality the annihilating operator
is over $\Q(x,y)$: $F$ has rational coefficients, and D-finiteness over $\C$
descends to the field of definition. Clearing denominators gives an operator of
some order $r$ with $x$-degrees at most some $D$. By
Theorem~\ref{thm:dichotomy}, $[\Q(\mu_H):\Q] \le D$ at all but at most $r$
heights $H \ge r$ --- contradicting Lemma~\ref{lem:northcott}, which says only
finitely many heights have degree $\le D$ while infinitely many heights exist.
\end{proof}

\section{The effective content}
\label{sec:effective}

Theorem~\ref{thm:main} is qualitative. Theorem~\ref{thm:dichotomy} is not: it
applies at a single height, so certified data at finitely many heights excludes
explicit boxes. This section gives two families of them, resting on different ingredients and
therefore independently usable.

\subsection{New-root content}

For the stored fixed-height generating functions $G_H = P_H/Q_H$ in lowest terms,
define the \emph{new-root content}
\[
  \psi_H \;:=\; \frac{Q_H}{\gcd\bigl(Q_H,\; Q_1Q_2\cdots Q_{H-1}\bigr)} ,
\]
the denominator factors appearing first at height $H$. Certified, for
$H = 1,\dots,10$:
\[
  \deg\psi_H \;=\; 1, 2, 4, 9, 29, 68, 181, 462, 1254, 3289 .
\]
Every $\psi_H$ is squarefree, and every $P_H/Q_H$ is in lowest terms, so every
root is active. The certificates are computed modulo the prime $p = 2^{61}-1$, where a
gcd equal to $1$ with preserved degrees is a rigorous certificate over
$\Q$.\footnote{The stored $H = 11$ entry is marked unvalidated and is excluded
here: its $Q_{11}$ shares no roots with $Q_9Q_{10}$ modulo $p$, which is
inconsistent with the structure at every validated level. It is a data-hygiene
flag, not a mathematical anomaly, and nothing above or below uses it. The stored
$Q_{11}$ is a CRT wraparound, wrong from degree $8159$ up.}

Run the argument of Theorem~\ref{thm:dichotomy} with \emph{all} roots of
$\psi_{H_0}$ in place of the single dominant one. Each is a pole of $G_{H_0}$ and,
by construction, a root of no earlier $Q_j$, so each of the lower slices in
\eqref{eq:rec} is a rational function finite there, and the
same argument gives $c(\alpha,H_0-m) = 0$ at $\deg\psi_{H_0}$ distinct values
$\alpha$. So $\deg\psi_{H_0} > D$ forces $c(\cdot,H_0-m) \equiv 0$ without any
appeal to the degree of $\mu_{H_0}$, and the rest of the proof is unchanged. The
rule is: exclude $(r,D)$ whenever at least $r+1$ certified levels $H \ge r$ have
$\deg\psi_H > D$.

\begin{table}[H]
\centering
\begin{tabular}{l l}
\toprule
order & excluded $x$-degrees \\
\midrule
$r \le 5$ & $D \le 28$ \\
$r \le 4$ & $D \le 67$ \\
$r \le 3$ & $D \le 180$ \\
$r \le 2$ & $D \le 461$ \\
$r \le 1$ & $D \le 1253$ \\
$r = 0$   & $D \le 3288$ \\
\bottomrule
\end{tabular}
\caption{No annihilating $y$-ODE for $F$ exists with these $(r,D)$, from the
$\psi$-degree certificates at $H \le 10$.}
\label{tab:psiboxes}
\end{table}

\subsection{Boxes from the degree of \texorpdfstring{$\mu_H$}{mu\_H}}

Both families read $\psi_H$ off the stored $G_H$; they differ in what they then
ask of it. The first extends Theorem~\ref{thm:dichotomy} to all roots of
$\psi_{H_0}$ at once, for which $\psi_H$ squarefree and $P_H/Q_H$ in lowest terms
is enough. The second uses Theorem~\ref{thm:dichotomy} exactly as proved, at the
dominant root only, and pays for that with irreducibility certificates,
which are what turn $\deg\psi_H$ into $[\Q(\mu_H):\Q]$. It is the weaker family
and is recorded because it exercises no more than the theorem above.

$\psi_H$ is irreducible over $\Q$ for $H = 1,\dots,8$ (degrees
$1, 2, 4, 9, 29, 68, 181, 462$) by single-prime certificates for $H \le 4$ and
multi-prime subset-sum-intersection certificates for $H = 5,\dots,8$. (At degrees
$1254$ and $3289$ the factoring cost was not paid.) Hence
$[\Q(\mu_H):\Q] = \deg\psi_H$ at every certified level.

\begin{table}[H]
\centering
\begin{tabular}{l l}
\toprule
order & excluded $x$-degrees \\
\midrule
$r \le 4$ & $D \le 8$ \\
$r \le 3$ & $D \le 28$ \\
$r \le 2$ & $D \le 67$ \\
$r \le 1$ & $D \le 180$ \\
$r = 0$   & $D \le 461$ \\
\bottomrule
\end{tabular}
\caption{Weaker than Table~\ref{tab:psiboxes}, and using
Theorem~\ref{thm:dichotomy} at the dominant root only.}
\label{tab:irrboxes}
\end{table}

\section{Transport to other families}
\label{sec:transport}

The proof of Theorem~\ref{thm:main} used exactly three things:

\begin{enumerate}
\item the counts $T(n,H)$ are nonnegative integers;
\item each height slice $G_H$ is rational, which follows from a finite column
  transfer matrix;
\item the slice growth constants $\mu_H$ are strictly increasing and bounded
  above.
\end{enumerate}

Nothing about king adjacency appears anywhere. So:

\begin{corollary}
\label{cor:transport}
The height-anisotropic generating function is not D-finite for fixed polyominoes
(the \oeis{A001168} triangle), for polyhexes, for polyiamonds, and generally for
any lattice-animal family admitting a column transfer matrix with a strongly
connected signature digraph and a finite growth constant.
\end{corollary}

The only per-family work is item~3, and both halves of it are cheap. Strictness
is the Perron--Frobenius argument of Lemma~\ref{lem:mono}, which needs the
signature digraph to be strongly connected --- true whenever a full column is a
legal state and every state can be closed off by one. Boundedness needs any crude upper bound on the family's growth constant, and a
Redelmeier-scan count of the kind in companion paper L3 always supplies one.

\section{Relation to existing work}
\label{sec:related}

\subsection{Bousquet-M\'elou and Rechnitzer 2002}

Their Lemma~9 is the criterion the haruspicy papers use, cited there
as ``(from~[4])'' and ``(from~[6])'' respectively; it is not in Haruspicy~1,
which supplies the combinatorial section and density machinery instead. Stated in
their notation: if $S(q,u) = \sum_n S_n(q)u^n$ is D-finite in $u$, and $P$ is the
union of the pole sets of the $S_n$, then $P$ has only finitely many limit points.

Side by side:

\begin{table}[H]
\centering\small
\begin{tabular}{p{0.20\textwidth} p{0.35\textwidth} p{0.35\textwidth}}
\toprule
 & \cite[Lemma~9]{bmr2002} & this paper (Thm.~\ref{thm:dichotomy}) \\
\midrule
opening move & extract $y$-coefficients of the ODE, giving a recurrence with polynomial coefficients & \textbf{identical}, and cited to them \\
what is tracked & the \emph{whole} pole set of every slice & \emph{one} pole per slice, $x = 1/\mu_H$ \\
finiteness input & topological: limit points of a subset of $\C$ & arithmetic: Northcott (bounded degree and bounded house) \\
mechanism & asymptotic in $n$: divide by $n^d$, take the limit, hit the leading coefficient & at a single level: lower faces regular at $1/\mu_{H_0}$, so $c_0$ vanishes there; degree comparison over $\Q$ \\
conclusion & qualitative & \textbf{effective}: explicit $(r,D)$ boxes \\
input needed & an explicit combinatorial description of the denominators, per family & Perron--Frobenius monotonicity plus any crude growth bound \\
\bottomrule
\end{tabular}
\caption{The shared opening, and where the two arguments part company.}
\label{tab:sidebyside}
\end{table}

Step~1 is theirs and dates to 2002. We cite it inside the proof rather than
presenting the derivation as self-contained. Our claim is that the remaining
hypotheses are light.

\subsection{Klazar 2003}

Klazar's Theorem~1: if $F' = G(x,F)$ with $G$ a Laurent series, $F$ \emph{not}
analytic, $G$ analytic and $\mathrm{ord}_y(G) < 0$, then $F$ is not $D_A$-finite
and a fortiori not D-finite. The lever is divergence, that $F$ must have zero
radius of convergence, and it cannot engage with our object at all: every
$G_H$ is rational with radius $1/\mu_H$, and $\mu_H < \lambda$, so $F$ is
analytic in a bidisc and his first hypothesis fails outright.

We record this because the negative is useful, a one-variable
non-$P$-recursiveness proof having looked as though it might be arithmetic, and
because Klazar has one thing we do not. He enlarges the target class from
D-finite to $D_A$-finite, with analytic rather than polynomial coefficients, on
the explicit ground that ``the key lies in the analytic $\times$ nonanalytic
dichotomy'', and so gets the stronger conclusion for free. Our theorem cannot follow: the degree comparison fails the instant the
coefficients are analytic, since an analytic function may vanish at
$1/\mu_{H_0}$ without vanishing identically. \textbf{Theorem~\ref{thm:main} is not-D-finite over $\Q(x)$ and does
not extend to $D_A$-finite.}

\subsection{Arithmetic neighbors}

Height theory has been applied to D-finiteness directly. Bell, Nguyen and
Zannier, \emph{D-finiteness, rationality, and height}~\cite{bnz2020}, with its
two sequels, is the nearest relative we found; Bell, Hu and
Satriano~\cite{bhs2020} is arithmetic dynamics along orbits of rational maps and
is further away. The configuration differs from ours in what is measured: they
bound the Weil heights of the \emph{coefficients} and conclude rationality; we
bound degree and house of the \emph{slice growth constants} and conclude a
contradiction. Bell, Gerhold, Klazar and Luca~\cite{bgkl2008} is a
zero-counting argument, again analytic.

We name these so that this section rests on the arithmetic neighborhood being
mapped, rather than on absence.

\subsection{What the searches establish}

Three independent databases were searched, and the record is per-search in
\repo{papers/MISSING.md}. A forward citation crawl over
\cite{bmr2002} (83 and 74 citing works by two indexes), the haruspicy
papers, \cite{chanRechnitzer2018}, \cite{bbo2021}, \cite{bhs2020} and
\cite{bgkl2008} found that every descendant proving non-D-finiteness uses pole or
singularity accumulation; no arithmetic endgame on a family of slice growth
constants appears in any of those sets. A third pass used full-text search within
citing articles --- a different instrument, and the one that could catch
a paper using Northcott inside a proof without saying so in an abstract: over
\cite{bmr2002}'s citing articles, ``Northcott'' returns zero hits and ``Mahler''
zero hits, while ``algebraic degree'' and ``height'' return hits that are all
false positives (the degree of an algebraic equation; column height in
bargraphs).

No collision found. Nothing in this paper is claimed on the strength of that:
we assert the lighter hypotheses and the effective boxes, and we cite
\cite[Lemma~9]{bmr2002} for the shared opening.

\section{Open problems}

\begin{openproblem}[the isotropic sequence]
Theorem~\ref{thm:main} says nothing about $a(x) = \sum_n a(n)x^n$, the row sums.
The same is true of Rechnitzer's SAP theorem, and for the same structural reason:
summing over $H$ destroys the slice structure the proof depends on. Nothing here
should be read as evidence about $a(x)$ either way.
\end{openproblem}

\begin{openproblem}[extend past $\Q(x)$]
Is $F$ non-$D_A$-finite? Our method cannot say, for the reason in
\S\ref{sec:related}. An argument that survives analytic coefficients would need
a different endgame than a degree comparison.
\end{openproblem}

\begin{openproblem}[certify the top two levels]
$\psi_9$ and $\psi_{10}$, at degrees $1254$ and $3289$, are squarefree and
new-root-certified but not certified irreducible. Doing so would upgrade
Table~\ref{tab:irrboxes} toward Table~\ref{tab:psiboxes}.
\end{openproblem}

\begin{openproblem}[quantitative Northcott]
Lemma~\ref{lem:northcott} is qualitative: for each $D$, finitely many heights
have degree $\le D$. An effective version, counting \emph{how many} $\mu_H$ can
have degree $\le D$ given house $\le 9.3154$, would turn the exclusion boxes from
data-dependent into proved, since the counting bound in the proof is explicit.
This looks like the most reachable of these four.
\end{openproblem}

\section{Reproduction}
\label{sec:repro}

\begin{table}[H]
\centering\small
\begin{tabular}{@{}l R{0.56\textwidth}@{}}
\toprule
claim & check \\
\midrule
$\psi$-degrees, squarefreeness, lowest terms & \repo{experiments/anisotropic\_dfinite.py} \\
Table~\ref{tab:psiboxes} and Table~\ref{tab:irrboxes} & \repo{experiments/anisotropic\_dfinite.py} \\
irreducibility certificates, $H \le 8$ & \repo{experiments/anisotropic\_dfinite.py} \\
Table~\ref{tab:mu} & \repo{experiments/mu\_H\_from\_atoms.py} \\
the stored $G_H$ & \repo{results/fixed\_height\_gfs.txt} \\
the house bound $\lambda \le 9.3154$ & companion paper; \repo{docs/proofs/polyplet-upper-bound.md} \\
the literature record & \repo{papers/MISSING.md}, \repo{experiments/citation\_crawl.py} \\
this paper's printed numbers & \repo{paper/verify\_l\_papers.py} \\
\bottomrule
\end{tabular}
\caption{Where each claim is checked. All certificates are computed modulo
$p = 2^{61}-1$; a gcd of $1$ with preserved degrees over $\F_p$ is a rigorous
certificate over $\Q$, since a factor over $\Q$ would reduce to a factor mod $p$.}
\label{tab:repro}
\end{table}

\bibliographystyle{plain}
\bibliography{shared/refs}

\end{document}
